This chapter presents the four empirical studies. The following sections detail each study’s motivation, empirical design, key findings, and contributions. While each study targets a specific decision and context, they collectively illustrate how data-driven approaches, grounded in consumer behavior theory, can be used to assess the effectiveness of strategies across formats, channels, and digital platforms.

\section{Study 1}

\textbf{From Price to Profit: Multi-Tier Private Label Response Across Retail Formats}

\noindent
\textit{Authorship:} Tanetpong Choungprayoon (Stockholm School of Economics)

\noindent
\textit{Status:} Working paper. An earlier draft was presented at the 53rd Annual Conference of the European Marketing Academy in 2024.

\noindent
***

Retailers increasingly use PLs not only to enhance margins but also to differentiate their assortments and target heterogeneous customer segments across store formats. Multi-tier PL strategies (typically comprising value, standard, and premium tiers) enable retailers to address heterogeneous customer preferences and strengthen store identity. However, the effectiveness of these tiers depends on store context. The same price change may have different consequences across formats that differ in assortment and in the shopping missions and customer expectations they facilitate. This study examines how price elasticities of multi-tier PLs (and the resulting profit effects of pricing decisions) vary across hypermarkets, supermarkets, and convenience stores.

Previous research has established the growing importance of PLs and the benefits of tiered strategies (\cite{Ailawadi2008-zy,Geyskens2010-bx}), but empirical evidence on how these effects differ across retail formats remains limited (\cite{Keller2022-tz}). While prior work shows that customer price responsiveness can vary across store formats (e.g.,\cite{Haans2011-cv}), existing studies often treat PL pricing as uniform, overlooking the possibility that store format moderates both customer price sensitivity and cross-tier interactions. This study addresses the research gap by quantifying how pricing responses differ across tiers and formats, and by linking these responses to profit outcomes.

To analyze the effectiveness of PL pricing across tiers and formats, the study employs a two-stage modeling framework linking demand responses to profitability outcomes. The specification captures the interdependence among PL tiers by incorporating cross-price effects across tiers. First, a sales response model is estimated using an error correction specification that includes both own-price and cross-price effects across PL tiers. This structure allows the model to capture within-category substitution and cannibalization effects that arise when price changes in (e.g.) the value tier influence sales in (e.g.) the premium tier. Short-term elasticities are normalized to enable comparison across categories and formats, and are used to simulate expected changes in sales volumes under different pricing conditions. Second, a profits-response model is estimated, also in error correction form. The model incorporates predicted changes in sales (derived from the sales response model) as key inputs, linking quantity adjustments back to pricing decisions. This step allows the estimation of format- and tier-specific profit elasticities, reflecting how price changes translate into bottom-line outcomes under real-world cannibalization and substitution dynamics. Both models include controls for seasonality, store-week effects, and endogeneity (via Gaussian copula terms), ensuring robustness of the estimates.

The findings reveal three key insights. First, customer price responsiveness varies significantly across formats, even within the same PL tier and category. Second, cross-tier pricing effects are nontrivial, meaning that offering a discount on a value tier product may boost volume for that tier but reduce sales in adjacent tiers, especially in formats with constrained assortments. Third, profitability outcomes are highly format-contingent. While premium PLs tend to be less elastic overall, their margin benefits depend on format-specific volume dynamics and price sensitivity.


\section{Study 2}
\textbf{Regular Prices vs. Discounts: How Pricing
Strategies Perform Across Store Formats}

\noindent
\textit{Authorship:} Tanetpong Choungprayoon (Stockholm School of Economics), Rickard Sandberg (Stockholm School of Economics), and Sara Rosengren (Stockholm School of Economics)

\noindent
\textit{Status:} Working paper. An earlier draft was presented at the 51st Annual Conference of the European Marketing Academy in 2022.

\noindent
***

Retailers frequently rely on two key pricing tactics to drive sales: adjustments to regular prices and the use of promotional discounts (i.e., temporary price cuts). While both tools can influence customer behavior, they are not equivalent in either perception or outcome. This study examines how regular price changes and promotional discounts differentially affect sales across retail formats within the same grocery chain. The analysis builds on the idea that shopping goals vary across formats, even for the same customers, yielding diverging responses to price cues.

Existing literature has highlighted the psychological distinctions between regular prices and promotional discounts (e.g., \cite{DelVecchio2006-xp,Shankar1996-kq}). While prior research has explored how customers respond to price changes, fewer studies have empirically disentangled the effects of regular price changes and promotional discounts across different store formats within the same retail chain. This study leverages an important feature of the empirical setting: promotional campaigns are largely centrally coordinated, while individual store managers retain flexibility in setting regular prices.

The analysis is based on data from a grocery retailer. It draws on three stores, one representing each format: a hypermarket, a supermarket, and a convenience store. All three are located in the same geographic area to minimize demographic and regional variation. Categories, products, and brands are selected based on consistency of purchase across format and time, allowing reliable cross-format comparison of price responsiveness. In the first stage of the analysis, we estimate regular price and discount elasticities for brand sales using an error-correction model of sales response following \textcite{Pauwels2007-ky}, with controls for category traffic and trip types (major, fill-in, and unplanned). These controls help account for differences in shopping behavior across formats. In the second stage, we regress the estimated elasticities on brand- and category-level attributes such as discount frequency, discount depth, PL status, storability, and category competitiveness to examine how these characteristics explain variation in pricing effectiveness across formats.

The findings reveal three key insights. First, regular price elasticities are negative across all store formats, with the strongest magnitude in hypermarkets and the weakest in convenience stores. Discount elasticities are also highest in hypermarkets, followed by supermarkets, indicating stronger responsiveness to promotional pricing in larger formats. Second, product characteristics matter: categories that are storable or impulsive exhibit higher discount sensitivity, while brands with deeper and more frequent promotions tend to experience greater discount elasticity. Third, depending on format and positioning, PLs show distinct patterns in their responses to price cues.



\section{Study 3}

\textbf{The Long-Term Effects of Online Channel Adoption in Grocery Retailing}
\noindent
\textit{Authorship:} Tanetpong Choungprayoon (Stockholm School of Economics), Emelie Fröberg (Stockholm School of Economics), and Sara Rosengren (Stockholm School of Economics))

\noindent
\textit{Status:} The current draft was initially given a “reject and resubmit” decision and was later rejected by the \textit{Journal of Retailing}.

\noindent
***

The introduction of digital channels has reshaped grocery retailing by offering consumers greater flexibility while posing new challenges for retailers. Many studies have examined who adopts online grocery shopping and why, but fewer have explored how purchasing behavior evolves in the long term after adoption. This study investigates how customer purchase behavior changes after the introduction of an online channel, and whether such adoption leads to sustained revenue growth for the retailer.

Previous research shows that adding new channels can benefit retailers because multichannel customers often spend more than single-channel shoppers (\cite{Ansari2008-qz,Kumar2005-nn}). Yet these effects may vary across product categories (\cite{Kushwaha2013-iu}) and depend on the specific benefits that the new channel provides (\cite{Bell2018-yn,Melis2016-mf}). In grocery retailing, habitual purchasing and the mix of hedonic and functional goods make long-term behavioral shifts particularly uncertain. Studies suggest that online channels enhance shopping utility for certain categories (e.g., bulky or heavy items that are convenient to order online) but reduce it for products requiring physical inspection or freshness evaluation (\cite{Campo2015-pt,Campo2020-xu}). However, less is known about whether these category-specific effects persist or attenuate over time.

To address this gap, our study builds on the theory of shopping utility maximization, which posits that consumers choose purchase patterns that yield the highest total utility. Shopping utility comprises two components: acquisition utility, referring to the benefits of the products offered (e.g., price and assortment), and transaction utility, referring to the convenience and effort associated with obtaining them (e.g., ordering and delivery). In this context, the retailer offered identical prices and assortments across channels, keeping acquisition utility constant. Differences in transaction utility therefore drive observed behavioral changes. The framework links three levels of outcomes: (a) category-level effects driven by perceived online ordering convenience, delivery convenience, and ordering risk; (b) customer-level effects on shopping frequency and monetary value; (c) retailer-level effects on total revenue.

The empirical analysis uses loyalty card data from a Nordic grocery retailer that launched an online channel in late 2015. The sample follows 578 households over a five-year period, including both adopters and non-adopters of the online channel. To identify long-term effects, the study applies a DiD design combined with PSM. It compares purchasing behavior before and after adoption while controlling for self-selection and customer heterogeneity.

The results show that the long-term effects of adoption are consistent with our theoretical framework. Customers who adopted the online channel spent more in product categories that offered high delivery convenience (e.g., bulky or heavy products) and spent less in sensory categories associated with higher perceived online purchase risk. In contrast, purchases of categories associated with online ordering convenience showed no significant long-term change. At the customer level, shopping frequency remained broadly stable while the average monetary value per purchase increased over time. These changes led to higher total expenditure among adopters, indicating higher long-term revenue for the retailer.

This study extends previous research on multichannel retailing by providing empirical evidence on the long-term effects of online channel adoption across category, customer, and retailer outcomes. It also illustrates how a framework based on shopping utility maximization can help retailers analyze influences of digital channel adoption on customer value over time.



\section{Study 4}
\textbf{How Did COVID-19 Affect Playlist Followers on Spotify?}

\noindent
\textit{Authorship:} Tanetpong Choungprayoon (Stockholm School of Economics), Hannes Datta (Tilburg University), and Max Pacheli (Tilburg University)

\noindent
\textit{Status:} Unpublished manuscript. An earlier draft was presented at Economics of the Music Industry in 2022.

\noindent
***

The COVID-19 pandemic reshaped how music was consumed and discovered, particularly on digital platforms where playlists have a central role in promoting track consumption. This study investigates how follower dynamics on Spotify playlists changed during the pandemic. The focus is on curator identity, playlist popularity, and share of major-label tracks. The purpose is to assess how an external disruption affected engagement and the distribution of visibility among stakeholders operating within a platform-controlled environment.

Previous research has shown that playlists have a central role in shaping music discovery and consumption on streaming platforms (\cite{Aguiar2018-wt,Datta2018-xz}). However, most existing evidence comes from periods of normal platform activity, and this limits understanding of how engagement patterns shift when external shocks disrupt user routines. Our study addresses the limit by examining how playlist follower growth responded to the pandemic. It offers insight on how platform exposure and curator outcomes evolved under disruption.

The study investigates these issues through four empirical questions related to overall follower changes, playlist popularity, curator type, and share of major-label tracks within playlists. Together, these questions assess whether engagement during the pandemic became (a) more concentrated among dominant stakeholders or (b) more broadly distributed across the platform.

The empirical analysis uses playlist-level data from Spotify, collected through third-party sources including Chartmetric for follower information and Every Noise at Once for playlist featuring indicators. The final sample, covering periods before and after the pandemic was declared, includes 39,918 active playlists tracked weekly from October 2019 to October 2020. Focusing on within-playlist variation, the empirical strategy employs a unit fixed effects model to estimate how playlist followers changed before and after the pandemic onset. This approach controls for persistent playlist characteristics (e.g., genre, theme, long-run popularity) while measuring the effect of time-varying pandemic indicators. Three complementary measures are used: a binary indicator for the pandemic declaration, a country-weighted government Stringency Index capturing the intensity of restrictions, and a country-weighted measure of residential mobility capturing changes in time spent at home.

The study also examines heterogeneous effects across playlists. Playlists are categorized into five popularity tiers based on follower percentiles, allowing the analysis to compare highly popular “superstar” playlists with mid-tail and long-tail playlists. In addition, playlists are classified according to curator type (including Spotify, major labels, independent labels, artists, professional curators, and users). Finally, the analysis considers the share of tracks from major record labels within each playlist to evaluate whether established industry content influenced follower growth during the pandemic.

The results show that there was no significant change in average follower growth following the pandemic declaration. However, an increase in follower growth when government restrictions intensified and users spent more time at home suggests that playlist engagement responded more to behavioral changes than to the declaration itself. The effects also differed across playlists. Mid-tail and long-tail playlists experienced relatively stronger growth compared to the most popular superstar playlists, indicating a broader distribution of engagement across content. Spotify-curated playlists were more resilient than those curated by other stakeholders, and (particularly among mid-tail and long-tail playlists) they showed stronger follower growth. Playlists with higher shares of major-label tracks showed some positive association with follower growth, but this effect was relatively small once track-level popularity was taken into account.

These findings demonstrate how large-scale disruptions can reshape engagement patterns on digital platforms. During the pandemic, follower growth became less concentrated among the most popular playlists, yet platform-curated playlists maintained a strong position within the ecosystem. The results highlight the role of platform governance, curator resources, and content composition in shaping visibility and engagement on streaming platforms across periods of external disruption.




